
Code generation models trained with single-objective optimization exhibit characteristic failure modes. Models optimized exclusively for execution feedback---learning from automated test case pass/fail signals---achieve functional correctness but produce unmaintainable code. Conversely, models aligned solely through human preference signals generate readable code that often fails execution tests. This tension forces practitioners to choose between code that works but is messy, or code that reads well but contains bugs.

Multi-objective reinforcement learning requires balancing competing reward structures. Execution feedback rewards any functional solution regardless of quality, while human feedback rewards elegant patterns regardless of correctness. Static integration approaches apply fixed weights throughout training, forcing a single compromise across all learning stages. Recent curriculum learning work schedules task difficulty, but no existing method explores curriculum over feedback modality---dynamically adjusting which feedback type dominates training phases.

We hypothesize that sequential capability building offers a solution. Just as developers first establish correctness before refining quality, then rely on experience for edge cases, training should emphasize execution feedback early, transition to AI-scaled quality feedback mid-training, and increase human oversight late in training. This represents a curriculum not over problem difficulty, but over feedback type.

This paper validates the mechanism of tri-modal reinforcement learning with dynamic feedback scheduling. Our framework integrates three heterogeneous reward signals---execution feedback from automated test cases, AI feedback from learned reward models, and human feedback from quality preferences---through a three-phase Gaussian schedule. We validate the mechanism through four sub-hypotheses testing predicted weight patterns and phase-specific objectives.

Our contributions are threefold. First, we design a tri-modal RL framework implementing dynamic feedback modality curriculum. Second, we validate the mechanism across four sub-hypotheses using competitive programming benchmarks, confirming predicted weight patterns emerge as designed. Third, we demonstrate each phase achieves its intended objective through gate criteria validation.

We disclose limitations transparently. Our experiments validate the training mechanism---weight scheduling logic, feedback collection, phase-specific objectives---but do not test end-to-end performance gains. Models used pretrained CodeGen-350M without actual reinforcement learning training. Human feedback uses heuristic-based quality proxies rather than real annotator ratings. These are acceptable for proof-of-concept mechanism validation, establishing foundation for future performance evaluation.

Our work opens a research direction for multi-objective RL: scheduling which feedback signal to emphasize based on training stage. The question shifts from ''which feedback signal?'' to ''which signal when?'' Mechanism validation provides evidence this approach is testable in full RL training.
